
In the case where we have a virtual object in the alias table $T$,
the approach puts it as the second component in the first cell of $T$. This behavior is dictated by our stack implementation.

% \noteAnt{Ce n'est pas très joli, mais ça évite de faire un IfThenElse qui est coûteux en temps,
% et qui entraine l'utilisation du prédicteur de branchement qui dans notre cas ca se tromper très souvent.
% Si on veut éviter le +1, il suffit de mettre le seuil en 3e position dans le triplet.}

% \noteAnt{dans la méthode random.choices de python, ce n'est pas l'alias qui est codé,
% pourtant au-delà d'une certain taille de simulation ça pourrait avoir du sens...
% faire des test}

Let us now introduce our sampling algorithm with virtual object rejection.
\begin{algorithm}[htb]
\caption{Random sampling with virtual object rejection} 
\label{algo:tableXY}
	\begin{algorithmic}[1]
        \scriptsize
        \renewcommand{\algorithmicrequire}{\textbf{Input:}}
        \renewcommand{\algorithmicensure}{\textbf{Output:}}
        \Require $T$, $\cs$, and $\texttt{virtual\_object}$.
        \Ensure A random object from $T \setminus \{\texttt{virtual\_object}\}$.
        \Function{RandomGenerationOld}{$T$, $\cs$, $\texttt{virtual\_object}$}
            \State $\bar{q} = \texttt{len}(T)$
            \State $result = \texttt{virtual\_object}$
            \While{$result == \texttt{virtual\_object}$}
                \State $i = \textsc{Uniform}(\bar{q})$
                \State $b = \textsc{Bernoulli}(T[i][0], \cs)$
                \If{b}
                    \State $result = T[i][1]$
                \Else
                    \State $result = T[i][2]$
                \EndIf
            \EndWhile
            \State \Return $result$
        \EndFunction
    \end{algorithmic}
\end{algorithm}

The way we implemented Algorithms $\textsc{Uniform}$ and $\textsc{Bernoulli}$ (called in Lines 5 and 6),
is related on Lumbroso's paper~\cite{Lumbroso13}.
He proved their optimality in terms of random bits, not of running time (\Cref{sec:exp} also evaluates a faster variant).
Even if both calls are optimal, 
the overall algorithm is not optimal in the random bits consumption.

\begin{proposition}
\label{prop:bits_consumption}
    The consumption of random bits is in average upper bounded by
    \[
        \begin{cases}
            \log_2(\bar{q}) + 3 \leq \log_2\!\big(n + \lfloor \frac{n-1}{\cs} \rfloor\big) + 3 & \text{if } r = 0,\\[1mm]
            \left(1+\frac{\cs - r}{w}\right) \left(\log_2(\bar{q}) + 3\right) \leq \left(1+\frac{1}{n}\right) \left(\log_2\!\big(n + \lceil \frac{n-1}{\cs} \rceil\big) + 3\right) & \text{if } r > 0.
        \end{cases}
    \]
    % In both cases, in particular, as soon as $\cs \geq n-1$, it is upper bounded by $\left(1+\frac{1}{n}\right) \left(\log_2(n+1) + 3\right)$.
\end{proposition}

We say that the sampled object is rejected $t$ times when the \textbf{While} loop 
(Line 4 in Algorithm~\ref{algo:tableXY}) is executed $t+1$ times.
\begin{lemma}
\label{lem:bits_consumption}
    Suppose $r > 0$. The probability that the sampled object is rejected $t$ times is
    \[
        \mathbb{P}(\text{exactly }t \text{ rejections}) = \left( \frac{\cs - r}{\bar{w}} \right)^t \left( \frac{w}{\bar{w}} \right).
    \]
    The average number of rejections is equal to $\displaystyle{\mathbb{E}(\text{number of rejections}) = \frac{\cs - r}{w}}$.
\end{lemma}
\begin{proof}
    The case when the sampled object is rejected has probability $\frac{\cs - r}{\bar{w}}$ that corresponds
    to the ratio of the weight of the virtual object and the total weight. 
    Thus having exactly $t$ rejections has probability $\left( \frac{\cs - r}{\bar{w}} \right)^t$
    and then the next object is valid (is not the virtual object) with probability $\left( \frac{w}{\bar{w}} \right)$.
    Thus the first statement is proved.
    To derive the average number of rejections, we have to calculate 
    \[
        \mathbb{E}(\text{number of rejections}) = \sum_{t=0}^{\infty} t \left( \frac{\cs - r}{\bar{w}} \right)^t \left( \frac{w}{\bar{w}} \right).  
    \]
    The next step is classical and can be derived from the derivative of the geometric series 
    $\sum_{t=1}^{\infty} t x^{t-1} = \frac{1}{(1-x)^2}$ for $|x|<1$.
    Hence,
    \begin{align*}
        \mathbb{E}(\text{number of rejections}) &= \frac{w}{\bar{w}} \cdot \frac{\cs - r}{\bar{w}} \cdot \sum_{t=1}^{\infty} t \left( \frac{\cs - r}{\bar{w}} \right)^{t-1} 
        %= \frac{w}{\bar{w}} \cdot \frac{\cs - r}{\bar{w}} \cdot \frac{1}{\left( 1 - \frac{\cs - r}{\bar{w}} \right)^2}
        %\\
        %    &= \frac{w}{\bar{w}} \cdot \frac{\cs - r}{\bar{w}} \cdot \frac{\bar{w}^2}{w^2} 
        = \frac{\cs - r}{w}.
    \end{align*}
    And the proof is complete. \qed
\end{proof}

The proof of \Cref{prop:bits_consumption} is a direct consequence of \Cref{lem:bits_consumption} and the fact that each iteration of the \textbf{While} loop performs one call to $\textsc{Uniform}(\bar{q})$ and one call to $\textsc{Bernoulli}$, which together consume on average at most $\log_2(\bar{q}) + 3$ random bits~\cite{Lumbroso13}.